\chapter{Sistemas de referencia y orientación}\label{ch:ejes}
\objetivos{Construir transformaciones pasivas; relacionar NED, cuerpo y viento; demostrar ortogonalidad; distinguir ángulos de actitud, aerodinámicos y de trayectoria.}

\section{Una convención completa}
Un vector físico no depende de los ejes elegidos, pero sus componentes sí. Se escribirá $\vect v^b=C_{bn}\vect v^n$: la matriz $C_{bn}$ transforma componentes desde el sistema $n$ al sistema $b$. El segundo subíndice identifica el origen de las componentes y el primero su destino. Es una transformación pasiva; no gira físicamente el vector.

El horizonte local $n$ tiene ejes norte, este y abajo, NED. El sistema cuerpo $b$ tiene $x_b$ hacia el morro, $y_b$ hacia el ala derecha y $z_b$ hacia abajo. Ambas ternas son dextrógiras. El momento positivo de alabeo baja el ala derecha, el momento positivo de cabeceo eleva el morro y el de guiñada gira el morro a la derecha.

El sistema inercial ideal tiene ejes sin rotación y origen sin aceleración. La Tierra gira respecto a él. El sistema terrestre ECEF gira con la Tierra; su eje $z$ se dirige al polo norte. NED depende de la posición: al desplazarse el avión sobre la esfera, el horizonte local gira. Para actuaciones locales se desprecia esa rotación y se aproxima NED como inercial. Para navegación extensa y para comparar con JSBSim habrá que recordar que el simulador incluye una Tierra rotatoria.

\section{De la Tierra al horizonte local}
Si $\lambda$ es latitud geocéntrica y $\ell$ longitud este, una transformación sobre Tierra esférica es
\begin{equation}
C_{ne}=\begin{bmatrix}
-\sin\lambda\cos\ell&-\sin\lambda\sin\ell&\cos\lambda\\
-\sin\ell&\cos\ell&0\\
-\cos\lambda\cos\ell&-\cos\lambda\sin\ell&-\sin\lambda
\end{bmatrix}.
\end{equation}
La primera fila apunta al norte y la tercera al centro terrestre. En una Tierra elipsoidal el abajo geodésico es normal al elipsoide, no radial; no se sustituye una latitud por otra sin revisar las definiciones. Para obtener coordenadas locales de una posición se resta primero la posición del origen y después se rota el vector diferencia. Rotar directamente dos posiciones con orígenes distintos no realiza una traslación.

\section{Tres rotaciones elementales}
Se adopta la secuencia intrínseca 3--2--1: guiñada $\psi$ sobre abajo local, cabeceo $\theta$ sobre el eje intermedio y alabeo $\phi$ sobre el eje longitudinal final. Las matrices pasivas son
\begin{align}
T_3(\psi)&=\begin{bmatrix}c_\psi&s_\psi&0\\-s_\psi&c_\psi&0\\0&0&1\end{bmatrix},\\
T_2(\theta)&=\begin{bmatrix}c_\theta&0&-s_\theta\\0&1&0\\s_\theta&0&c_\theta\end{bmatrix},\\
T_1(\phi)&=\begin{bmatrix}1&0&0\\0&c_\phi&s_\phi\\0&-s_\phi&c_\phi\end{bmatrix},
\end{align}
donde $c_\phi=\cos\phi$ y $s_\phi=\sin\phi$, y análogamente para los otros ángulos. La primera transformación que actúa sobre el vector aparece a la derecha. El producto es
\begin{equation}
C_{bn}=T_1T_2T_3=
\begin{bmatrix}
c_\theta c_\psi&c_\theta s_\psi&-s_\theta\\
s_\phi s_\theta c_\psi-c_\phi s_\psi&s_\phi s_\theta s_\psi+c_\phi c_\psi&s_\phi c_\theta\\
c_\phi s_\theta c_\psi+s_\phi s_\psi&c_\phi s_\theta s_\psi-s_\phi c_\psi&c_\phi c_\theta
\end{bmatrix}.\label{eq:Cbn}
\end{equation}
Una secuencia diferente produce otra matriz. Para comprobarlo basta girar $90^\circ$ alrededor de dos ejes y observar a dónde llega el tercero: los productos no conmutan.

\subsection{Pruebas geométricas}
Cada fila contiene los cosenos directores de un eje cuerpo respecto a NED. Las filas tienen norma unidad y son ortogonales. Para cada rotación elemental, $T_iT_i^T=I$ y $\det T_i=1$. Por tanto,
\begin{equation}
C_{bn}C_{bn}^{T}=T_1T_2T_3T_3^TT_2^TT_1^T=I,
\qquad \det C_{bn}=1.
\end{equation}
La inversa es $C_{nb}=C_{bn}^T$. En consecuencia, se conserva $\vect v\cdot\vect v$. Una matriz con determinante $-1$ puede conservar normas, pero incluye una reflexión y no describe sólo una rotación de una terna dextrógira.

\caso{Ejemplo: gravedad en actitud de cabeceo}{Para $\phi=0$, $\theta=10^\circ$, la gravedad local $\vect g^n=(0,0,g)^T$ tiene componentes $\vect g^b=(-g\sin10^\circ,0,g\cos10^\circ)^T=(-1.703,0,9.658)\SIu{m\,s^{-2}}$. El morro alto hace que la gravedad frene en $x_b$. La norma sigue siendo $9.80665\SIu{m\,s^{-2}}$. Si el primer componente resultase positivo, habría un error de signo o se estaría usando la transformación inversa.}

\section{Velocidad angular y ángulos de Euler}
Las derivadas de los ángulos no son las componentes de velocidad angular, porque las tres rotaciones se producen alrededor de ejes diferentes. Sumando sus vectores expresados en cuerpo se obtiene
\begin{equation}
\begin{bmatrix}p\\q\\r\end{bmatrix}=
\begin{bmatrix}1&0&-\sin\theta\\0&\cos\phi&\sin\phi\cos\theta\\0&-\sin\phi&\cos\phi\cos\theta\end{bmatrix}
\begin{bmatrix}\dot\phi\\\dot\theta\\\dot\psi\end{bmatrix}.
\label{eq:Euler-directa}
\end{equation}
El determinante es $\cos\theta$. Cuando no es nulo,
\begin{align}
\dot\phi&=p+(q\sin\phi+r\cos\phi)\tan\theta,\\
\dot\theta&=q\cos\phi-r\sin\phi,\\
\dot\psi&=(q\sin\phi+r\cos\phi)/\cos\theta.
\end{align}
En $\theta=\pm90^\circ$ se pierde unicidad en la parametrización. El avión no tiene una singularidad física; la tienen estas coordenadas. Cerca de ella, pequeñas incertidumbres en actitud producen grandes variaciones de guiñada y alabeo. Una integración robusta debe cambiar de representación o usar cuaterniones.

\section{Ejes viento y ángulos aerodinámicos}
Los ángulos aerodinámicos se construyen con la velocidad respecto al aire, no con la velocidad respecto al suelo. Sean $(u_a,v_a,w_a)$ sus componentes cuerpo y $V=\sqrt{u_a^2+v_a^2+w_a^2}>0$:
\begin{equation}
\alpha=\operatorname{atan2}(w_a,u_a),\qquad \beta=\arcsin(v_a/V).
\end{equation}
Se escoge $x_w$ paralelo a esa velocidad y $z_w$ abajo en ausencia de resbalamiento. La transformación correspondiente es
\begin{equation}
C_{bw}=\begin{bmatrix}
\cos\alpha\cos\beta&-\cos\alpha\sin\beta&-\sin\alpha\\
\sin\beta&\cos\beta&0\\
\sin\alpha\cos\beta&-\sin\alpha\sin\beta&\cos\alpha
\end{bmatrix}.\label{eq:Cbw}
\end{equation}
Su primera columna reproduce $\vect v_a^b/V$. La fuerza aerodinámica se escribe
\begin{equation}
\vect F_a^w=(-D,Y_w,-L)^T,\qquad \vect F_a^b=C_{bw}\vect F_a^w.
\end{equation}
Con $\beta=0$, $X_a=-D\cos\alpha+L\sin\alpha$ y $Z_a=-D\sin\alpha-L\cos\alpha$. Una fuerza de sustentación positiva no coincide, por tanto, con un $Z_b$ positivo. Los ejes de estabilidad se obtienen girando sólo por la incidencia de referencia; no coinciden con ejes viento instantáneos cuando hay perturbaciones.

\section{Actitud, rumbo y trayectoria}
El cabeceo describe la orientación del fuselaje respecto al horizonte. La trayectoria describe la dirección de la velocidad. En vuelo simétrico sin viento y con alas niveladas, $\theta=\alpha+\gamma$. En presencia de alabeo, resbalamiento o viento la suma escalar deja de ser general y deben emplearse matrices. El rumbo del morro $\psi$ tampoco tiene por qué coincidir con la dirección de desplazamiento sobre el terreno.

\section{Ampliación: cuaterniones}
Un cuaternión unidad puede escribirse $Q=(q_0,\vect q)$, con parte escalar primero y producto de Hamilton. Si representa la rotación activa de cuerpo a NED, y $\vect\omega^b$ es la velocidad angular de cuerpo respecto a un NED considerado inercial,
\begin{equation}
\dot Q=\tfrac12 Q\otimes(0,\vect\omega^b),\qquad q_0^2+\vect q^T\vect q=1.
\end{equation}
La matriz activa asociada es
\begin{equation}
C_{nb}=(q_0^2-\vect q^T\vect q)I+2\vect q\vect q^T+2q_0[\vect q]_\times.
\end{equation}
Su transpuesta transforma componentes NED a cuerpo. Las fórmulas cambian de signo o de orden al elegir otra convención de cuaterniones; declarar sólo ``uso cuaterniones'' no basta. $Q$ y $-Q$ representan la misma actitud. Normalizar periódicamente limita el error de norma, pero no corrige un signo equivocado en la ecuación cinemática.

\verifica{Las matrices Euler y viento han sido contrastadas simbólicamente en ortogonalidad y dirección de la velocidad; también se ensayaron 200 transformaciones numéricas de ida y vuelta. La representación de cuaterniones se presenta como ampliación con su convención explícita; no se identifica automáticamente con el orden interno de cualquier programa.}
\section{Taller resuelto: Una transformación completa de ida y vuelta}

\noindent\textbf{Nivel:} Intermedio.

\apartado{1. Enunciado}{Construya Cbn para alabeo 18 grados, cabeceo 9 y guiñada 32. Transforme el vector NED (96,8,-5) m/s a cuerpo, vuelva a NED y compruebe norma y determinante.}

\apartado{2. Datos}{Vector aéreo y ángulos Euler 3--2--1; no hay traslación del origen porque se transforma una velocidad.}

\apartado{3. Incógnitas}{Las magnitudes y comprobaciones solicitadas en el caso.}

\apartado{4. Sistema de referencia}{Matriz pasiva NED a cuerpo, vectores columna.}

\apartado{5. Hipótesis físicas}{Horizonte fijo local y ángulos fuera de singularidad de Euler.}

\apartado{6. Modelo matemático}{Una transformación completa de ida y vuelta}

\apartado{7. Ecuaciones simbólicas}{Se parte de los balances y definiciones siguientes.}

\begin{equation}C_{bn}=T_1(\phi)T_2(\theta)T_3(\psi)\end{equation}

\begin{equation}\vect v^b=C_{bn}\vect v^n,\quad\vect v^n=C_{bn}^T\vect v^b\end{equation}

\apartado{8. Desarrollo algebraico}{Las operaciones se realizan en este orden, conservando las condiciones del modelo.}

\begin{enumerate}

\item Se forman las tres matrices elementales con los signos pasivos.

\item Se multiplica en orden, empezando por guiñada a la derecha.

\item Se verifica la inversa por transposición y se compara con las componentes iniciales.

\item La norma se obtiene antes y después de la rotación y el determinante se calcula independientemente.

\end{enumerate}

\apartado{9. Sustitución numérica}{Se usan las magnitudes del apartado de datos en SI; los ángulos se convierten a radianes antes de evaluar funciones trigonométricas. Las densidades indicadas como ISA se obtienen de la capa correspondiente.}

\noindent u cuerpo:\quad $85.37962577153047\simeq 85.37963\,\mathrm{m\,s^{-1}}$.\par

\noindent v cuerpo:\quad $-39.31561924237238\simeq -39.31562\,\mathrm{m\,s^{-1}}$.\par

\noindent w cuerpo:\quad $21.67029272303927\simeq 21.67029\,\mathrm{m\,s^{-1}}$.\par

\apartado{10. Resultados con unidades}{Los resultados pertenecen al modelo y al intervalo de datos declarados.}

\begin{center}\begin{tabular}{lr}\toprule Magnitud&Resultado\\\midrule

u cuerpo & $85.37963\,\mathrm{m\,s^{-1}}$\\

v cuerpo & $-39.31562\,\mathrm{m\,s^{-1}}$\\

w cuerpo & $21.67029\,\mathrm{m\,s^{-1}}$\\

Norma & $96.46243\,\mathrm{m\,s^{-1}}$\\

Determinante & $1\,$\\

\bottomrule\end{tabular}\end{center}

\apartado{11. Interpretación}{La misma velocidad tiene componentes diferentes; su módulo y el vector físico se conservan.}

\apartado{12. Comprobación independiente y alcance}{La multiplicación numérica se reproduce en el programa y se comprueba C transpuesta por C e ida y vuelta.}

\apartado{13. Limitaciones}{Una matriz ortogonal con otro orden de ángulos podría conservar norma y aun representar una actitud equivocada; la prueba geométrica de secuencia sigue siendo necesaria.}

\par\smallskip\noindent{\footnotesize\textit{Datos y solución originales. Las sustituciones registradas se reproducen con SymPy y Studium; eso no establece validación experimental ni revisión experta.}}


\repaso{Con $\alpha=8^\circ$, $\beta=0$ y $V=70\SIu{m\,s^{-1}}$, ¿es $w_a$ negativo porque el avión asciende?}{No. $w_a=V\sin\alpha=9.74\SIu{m\,s^{-1}}$ es positivo en cuerpo. La velocidad vertical sobre el terreno depende además de la actitud y del viento.}
\resumen{Una transformación necesita ejes, sentido, orden y referencia de los ángulos. Las rotaciones propias conservan normas y tienen determinante unidad. Los ángulos de Euler parametrizan la actitud con singularidades; los aerodinámicos describen la velocidad relativa al aire.}
\fuentes{How, MIT 16.333, lección 3, pp. internas 3--1 a 3--9 (Euler y cuaterniones); lección 4, pp. 4--2 a 4--4 (ejes). Las matrices se derivan y comprueban en el proyecto. La notación se declara de nuevo para evitar trasladar signos de otros convenios.}
